\subsection{Phát biểu bài toán}
Trong việc quản lý gia phả của một dòng họ, thách thức lớn nhất nằm ở tính phân cấp và sự biến động của dữ liệu. Cụ thể, bài toán đặt ra các yêu cầu và ràng buộc sau:
\begin{itemize}
	\item Thứ nhất, một cá nhân có thể là con của người này nhưng lại là cha/mẹ của người khác. Mối quan hệ này kéo dài qua nhiều thế hệ, tạo thành một cấu trúc cây.
	\item Thứ hai, khác với cây nhị phân (mỗi nút chỉ có tối đa 2 con), một người trong thực tế có thể có số lượng con bất kỳ. Việc sử dụng mảng cố định để lưu trữ danh sách con sẽ gây lãng phí bộ nhớ hoặc thiếu linh hoạt.
	\item Thứ ba, cần lưu trữ rõ ràng các thông tin thực thể (Tên, giới tính, ngày sinh...) và các mối quan hệ (Cha-con, Anh-chị-em).
	\item Cuối cùng, bài toán yêu cầu các thao tác: thêm mới, duyệt và tìm kiếm.
\end{itemize}
\subsection{Cấu trúc dữ liệu}
\subsubsection{Thiết kế chung}

Để giải quyết bài toán quản lý thông tin đa dạng của mỗi cá nhân trong dòng họ, chương trình thiết kế một kiểu dữ liệu tự định nghĩa (\texttt{struct person}). Sơ đồ cấu trúc được mô tả ở Hình \ref{fig:ERD1}. Cấu trúc này ánh xạ trực tiếp các yêu cầu của thực tế vào không gian bộ nhớ:
\begin{itemize}
	\item Nhóm thuộc tính định danh và nhân khẩu học: Các trường dữ liệu \texttt{name}, \texttt{gender}, \texttt{birthday}, \texttt{job} được sử dụng để lưu trữ hồ sơ cơ bản. Trong đó, \texttt{gender} (Nam/Nữ) đóng vai trò là cờ điều hướng (flag) để quyết định cấu trúc phát triển của cây; trường \texttt{birthday} đóng vai trò định hình thứ tự anh - em trong những người con của cùng một cha.
	\item Nhóm thuộc tính gia đình và tình trạng: Trường \texttt{spouseName} (Tên vợ/chồng) được gắn trực tiếp như một thuộc tính của Nút thay vì tạo một Nút độc lập, giúp giảm độ phức tạp của đồ thị cây. Ngoài ra, kỹ thuật này đảm bảo cấu trúc gia phả luôn là một cây đơn gốc, giúp duy trì tính toàn vẹn cho các thuật toán đệ quy và tìm tổ tiên chung (LCA) mà không làm tăng độ phức tạp của đồ thị thành dạng chu trình. Trường deathDay quản lý tình trạng sinh tử, và \texttt{numChildren} (số con) được thiết kế đặc biệt chỉ kích hoạt lưu trữ đối với thành viên Nữ. 
	\item Nhóm thuộc tính định tuyến (Routing Attributes): Các con trỏ \texttt{pParent}, \texttt{pFir-stChild}, \texttt{pNextSibling} đóng vai trò là ``chỉ dẫn đường" để móc nối thực thể này vào mạng lưới gia phả chung, đảm bảo tính liên kết động mà không cần mảng trung gian.
\end{itemize}
\begin{figure}[!h]
	\centering
    \begin{tikzpicture}[
    node distance=2.5cm,
    transform shape,
    % Tự định nghĩa lại các khối đồ họa cơ bản nhất
    entity/.style={
        rectangle, draw=blue!80, fill=blue!10, 
        very thick, minimum width=3.5cm, minimum height=1.2cm, 
        font=\bfseries, drop shadow
    },
    attribute/.style={
        ellipse, draw=orange!80, fill=yellow!10, 
        thick, minimum width=2.5cm, minimum height=0.8cm, 
        font=\small, drop shadow
    },
    relationship/.style={
        diamond, draw=gray!80, fill=gray!10, 
        thick, aspect=1.5, text width=2.5cm, align=center,
        font=\small\bfseries, drop shadow
    },
    >={Stealth[round]},
    scale=0.8
]

    % 1. THỰC THỂ TRUNG TÂM
    \node[entity] (person) {PERSON (NODE)};

    % 2. CÁC THUỘC TÍNH (Gộp bớt để sơ đồ thoáng)
    \node[attribute] (name) [above=1.5cm of person] {\underline{Name}} edge[thick] (person);
    \node[attribute] (info) [above right=1cm and 1cm of person] {Profile Info} edge[thick] (person);
    \node[attribute] (id) [above left=1cm and 1cm of person] {ID (Tạm)} edge[thick] (person);

    % 3. QUAN HỆ FIRST CHILD (Màu đỏ, vòng về bên trái)
    \node[relationship, draw=red!70, fill=red!5] (rc) [below left=2cm and 0.5cm of person] {has\\FirstChild};
    \draw[-, thick] (person.210) -- (rc.north);
    \draw[->, red!70, thick] (rc.west) -- ++(-1.5,0) |- (person.190) 
        node[pos=0.75, above, font=\bfseries, text=red!70!black] {pFirstChild (1:1)};

    % 4. QUAN HỆ NEXT SIBLING (Màu cam, vòng về bên phải)
    \node[relationship, draw=orange!70, fill=orange!5] (rs) [below right=2cm and 0.5cm of person] {has\\NextSibling};
    \draw[-, thick] (person.330) -- (rs.north);
    \draw[->, orange!70, thick] (rs.east) -- ++(1.5,0) |- (person.350) 
        node[pos=0.75, above, font=\bfseries, text=orange!70!black] {pNextSibling (1:1)};

    % 5. QUAN HỆ PARENT (Màu xanh, vòng cung rộng ôm trọn bên trái)
    \node[relationship, draw=green!60!black, fill=green!5] (rp) [below=3.5cm of person] {has\\Parent};
    \draw[-, thick] (person.270) -- (rp.north);
    % Vẽ đường Parent vòng xa ra ngoài để không đè lên FirstChild
    \draw[->, green!60!black, thick] (rp.south) -- ++(0,-0.8) -| ([xshift=-4.5cm]person.west) |- (person.165) 
        node[pos=0.75, above, font=\bfseries, text=green!60!black] {pParent (N:1)};

\end{tikzpicture}
\caption{Sơ đồ ERD mô tả thực thể Person và các mối quan hệ tự thân trong cấu trúc cây LCRS}
\label{fig:ERD1}
\end{figure}



\subsubsection{Cài đặt các ràng buộc cần thiết}
Cấu trúc dữ liệu đã được tùy biến để tuân thủ nghiêm ngặt các quy tắc truyền thống của một gia phả dòng họ, gồm hai ý chính:
\begin{enumerate}
	\item Theo yêu cầu đề tài, con gái khi lập gia đình sẽ tính theo gia phả nhà chồng. Do đó, trong giải thuật \texttt{AddChild} (Thêm con), nhóm đã cài đặt chốt chặn logic, nghĩa là: Khi người dùng thực hiện thao tác thêm thế hệ mới, chương trình sẽ kiểm tra \texttt{parent->gender}. Nếu giá trị là ``Nữ", thao tác bị từ chối và nút đó vĩnh viễn là một ``Nút lá" (không bao giờ có \texttt{pFirstChild}).
	\item Trong một dòng họ, việc trùng họ tên giữa các thế hệ là khá phổ biến (ví dụ: nhiều người cùng tên ``Nguyễn Văn A"). Việc tìm kiếm và cập nhật của chương trình nếu chỉ dựa vào chuỗi name sẽ dẫn đến sai lệch nghiêm trọng. Để giải quyết, nhóm cơ chế Lọc danh sách lựa chọn (Selection List). Thay vì định danh bằng chuỗi ký tự không duy nhất, chương trình sử dụng địa chỉ bộ nhớ làm định danh tuyệt đối, kết hợp với việc truy vấn thuộc tính phụ từ người dùng để xác nhận đối tượng.
\end{enumerate}

\subsection{Thuật toán và đánh giá thuật toán}
\subsubsection{Thuật toán Duyệt và Tìm kiếm}

Vì các nút bộ nhớ (\texttt{Person}) được cấp phát động trên RAM, chương trình thiết kế hai chiến lược truy vấn chuyên biệt dựa trên loại khóa tìm kiếm để tối ưu hóa hiệu năng hệ thống:

\textbf{Truy vấn theo Định danh (ID Lookup)}
	\begin{itemize}
		\item \textit{Cơ sở lý thuyết:} Sử dụng cấu trúc Bảng băm (Hash Table) nhằm thiết lập ánh xạ trực tiếp $f: ID \rightarrow \text{Memory Address}$, cho phép truy xuất dữ liệu độc lập với độ sâu của đồ thị cây.
		\item \textit{Phương pháp thực hiện:} Kỹ thuật này được kích hoạt trong tiến trình đọc tệp nhị phân \texttt{.dat}. Chương trình khởi tạo \texttt{std::unordered\_map<int, Person*>} để lưu trữ danh sách các nút vừa cấp phát. Khi cần khôi phục tập hợp cạnh $E$ (nối con trỏ cha-con), hệ thống chỉ cần truyền \texttt{parentId} vào Bảng băm để nhận lại con trỏ thực thi ngay lập tức. Hình \ref{fig:Hash1} mô tả dễ hiểu quá trình này.
		\item \textit{Phân tích độ phức tạp:} Thuật toán đưa độ phức tạp thời gian tìm kiếm cha từ $O(|V|)$ xuống $O(1)$ tuyệt đối, giải quyết triệt để nút thắt cổ chai khi hệ thống phải tái thiết lập hàng vạn liên kết con trỏ.
	\end{itemize}

\vspace{0.5cm}
	
\begin{figure}[h!]
	\vspace*{-\baselineskip}
	\centering
	\begin{adjustbox}{center=-3cm}
	\begin{tikzpicture}[
		>=stealth, thick,
		box/.style={draw=black, fill=white, rounded corners, inner sep=10pt, align=left},
		highlight/.style={draw=blue!80, fill=blue!10, thick},
		arrowtext/.style={midway, above, font=\scriptsize\itshape, text=blue!80!black},
		transform shape,
		scale=0.65
		]
		
		\node[box, highlight] (file) {
			\textbf{Bản ghi đang đọc (File .dat)}\\
			\texttt{id:} 2 \\
			\texttt{parentId:} \textcolor{red}{\textbf{1}} \\
			\texttt{name:} Nguyễn Văn Hùng
		};
		
		% Khối 2: Bảng Băm (Hash Table)
		\node[box, right=3.5cm of file] (hash) {
			\textbf{Bảng Băm (\texttt{unordered\_map})}\\
			Khóa (ID) $\rightarrow$ Giá trị (Địa chỉ RAM)\\
			\textcolor{red}{\textbf{1}} \hspace{0.8cm} $\rightarrow$ \textcolor{green!50!black}{\texttt{0x1A2B}} \\
			2 \hspace{0.8cm} $\rightarrow$ \texttt{0x3C4D} \\
			... \hspace{0.6cm} $\rightarrow$ ...
		};
		
		% Khối 3: Cây LCRS trên RAM
		\node[box, right=3.5cm of hash] (ram_parent) {
			\textbf{Node Cha (Trên RAM)}\\
			\texttt{Địa chỉ: }\textcolor{green!50!black}{\texttt{0x1A2B}} \\
			\texttt{name:} Nguyễn Văn Tèo \\
			\texttt{pFirstChild:} \textcolor{orange}{\texttt{nullptr}} $\rightarrow$ \texttt{0x3C4D}
		};
		
		\node[box, below=1.5cm of ram_parent] (ram_child) {
			\textbf{Node Con mới tạo}\\
			\texttt{Địa chỉ: 0x3C4D} \\
			\texttt{name:} Nguyễn Văn Hùng
		};
		
		% Vẽ mũi tên liên kết
		\draw[->, red, very thick] (file.east) -- (hash.west) node[arrowtext, font=\small] {1. Truy vấn parentId};
		\draw[->, green!50!black, very thick] (hash.east) -- (ram_parent.west) node[arrowtext, font=\small] {2. Trả về địa chỉ};
		\draw[->, orange, very thick] (ram_parent.south) -- (ram_child.north) node[arrowtext, right, font=\small] {3. Nối con trỏ nhánh};
		
	\end{tikzpicture}
\end{adjustbox}
	\caption{Mô phỏng quy trình thực hiện truy vấn theo Định danh bằng cấu trúc Bảng băm}
	\label{fig:Hash1}
\end{figure}


\begin{figure}[H]
	\vspace*{-\baselineskip}
	\centering
	\begin{minipage}{0.50\textwidth}
		\begin{algorithm}[H]
			\footnotesize
			\caption*{Pha 1: Cấp phát nút}
			\begin{algorithmic}
				\State Khởi tạo bảng băm: HashMap<int, Person*>
				\For {mỗi Record trong fileRecords}
				\State pNew = New Person(Record.name, Record.gender)
				\State Sao chép thông tin từ Record sang pNew
				\State HashMap[Record.id] = pNew
				\EndFor
			\end{algorithmic}
		\end{algorithm}
	\end{minipage}%
	\hfill
	\begin{minipage}{0.50\textwidth}
		\begin{algorithm}[H]
			\footnotesize
			\caption*{Pha 2: Thiết lập liên kết}
			\begin{algorithmic}
				\For {mỗi Record trong fileRecords}
				\State currentPerson = HashMap[Record.id]
				\If {Record.parentId != 0}
				\State Tra cứu địa chỉ Cha từ Bảng băm
				\State Thực hiện nối nút
				\EndIf
				\EndFor
				
				\State Cập nhật global\_id\_counter = Max(Record.id) + 1
				\State Trả về địa chỉ Gốc (Root)
			\end{algorithmic}
		\end{algorithm}
		
	\end{minipage}      
	\caption{Mã giả quy trình thực hiện truy vấn theo Định danh bằng cấu trúc Bảng băm}
\end{figure}
	
\textbf{Truy vấn theo Tên gọi}
	\begin{itemize}
		\item \textit{Cơ sở lý thuyết:} Do thuộc tính ``Tên'' là dữ liệu chuỗi phi tuyến tính (không tạo thành cây tìm kiếm nhị phân BST), hệ thống triển khai thuật toán Duyệt theo chiều sâu (Depth-First Search - DFS).
		\item \textit{Phương pháp thực hiện:} Thuật toán sử dụng đệ quy Tiền thứ tự (Pre-order), duyệt tuần tự qua các liên kết \texttt{pFirstChild} và \texttt{pNextSibling}. Để khắc phục điểm yếu của thuật toán tham lam (Greedy) truyền thống, chương trình không trả về kết quả khớp đầu tiên. Thay vào đó, thuật toán tiếp tục duyệt cạn không gian đỉnh để thu thập toàn bộ các thực thể trùng tên vào một \texttt{std::vector<Person*>}. Khối kết quả này sau đó được đẩy lên tầng giao diện (UI) để người dùng phân biệt và chọn lựa mục tiêu chính xác.
		\item \textit{Phân tích độ phức tạp:} Độ phức tạp thời gian là $O(|V|)$ (số lượng thành viên). Độ phức tạp không gian lưu trữ cho ngăn xếp gọi hàm (Call Stack) là $O(H)$, với $H$ là chiều cao cực đại của cây gia phả. Với quy mô dữ liệu thực tế của một dòng họ, thuật toán đảm bảo tốc độ phản hồi ở mức thời gian thực.
	\end{itemize}

\subsubsection{Thuật toán Xác định quan hệ}
\label{sssec:lca1}

Bài toán Xác định quan hệ và danh xưng trong gia phả thực chất là một bài toán kinh điển của cấu trúc dữ liệu cây. Dưới góc độ cơ sở toán học trong khoa học máy tính, đây chính là sự kết hợp giữa thuật toán Tìm tổ tiên chung gần nhất (\textbf{L}owest \textbf{C}ommon \textbf{A}ncestor - viết tắt là LCA) và việc phân tích Ma trận độ lệch khoảng cách.

Để đảm bảo thuật toán chạy chính xác, cũng như bao quát được sự phức tạp của hệ thống danh xưng trong tiếng Việt, thuật toán sẽ trải qua bốn pha như sau:

\textbf{\textit{Pha thứ nhất: Truy vết đường đi về gốc.}}
\begin{itemize}
    \item Mục tiêu của pha: Tìm chuỗi các đỉnh kết nối từ người đang xét lên tới Cụ Tổ, làm cơ sở cho việc xác định LCA ở pha thứ hai.
    \item Phương pháp thực hiện: Giả sử đồ thị cây gia phả là $T = (V, E)$, ta cần xây dựng hai tập hợp đường đi cho nút $A$ và nút $B$. Quy trình thực hiện như sau:
    \begin{itemize}
        \item Từ con trỏ của $A$, dùng vòng lặp di chuyển theo hướng con trỏ {\tt pParent} cho đến khi gặp {\tt nullptr} (Nút gốc).
        \item Lưu lại các địa chỉ bộ nhớ đã đi qua vào một mảng (hoặc \texttt{std::vector}). Ta thu được tập $P_A = \{A, P_{A1}, P_{A2}, \dots, \text{Root}\}$.
        \item Làm tương tự với $B$, ta thu được tập $P_B = \{B, P_{B1}, P_{B2}, \dots, \text{Root}\}$.
    \end{itemize}
    \item Về độ phức tạp: Độ phức tạp thời gian xấp xỉ $O(H)$, với với $H$ là chiều cao lớn nhất của cây (số lượng thế hệ). Thao tác này cực kỳ nhanh vì nó chỉ đi theo một đường thẳng lên trên, không cần duyệt qua các nhánh khác của gia phả; Độ phức tạp không gian là $O(H)$ vì cần cấp phát hai mảng để lưu trữ tối đa $H$ con trỏ địa chỉ bộ nhớ cho mỗi đường đi.
    \end{itemize}

\begin{figure}[H]
	\vspace*{-\baselineskip}
    \centering
    \begin{minipage}{0.50\textwidth}
        \begin{algorithm}[H]
        	\footnotesize 
            \caption{GetPathToRoot(Node)} \label{alg:getpathtoroot}
            \begin{algorithmic}
                \State Khởi tạo mảng Path rỗng
                \While{Node không phải là rỗng}
                    \State Thêm Node vào đầu mảng Path (để Gốc luôn ở vị trí số 0)
                    \State Node = Node.parent
                \EndWhile
                \State \Return Path
            \end{algorithmic}
        \end{algorithm}
    \end{minipage}%
    \hfill
    \begin{minipage}{0.50\textwidth}
        \centering
        \begin{tikzpicture}[level distance=1.5cm, sibling distance=2.5cm]
        	% Tô màu toàn bộ đường đi từ A lên Root và B lên Root
        	\node (root) [person, pathA, pathB] {Cụ Tổ}
        	child { node (c1) [person, pathA] {Cha A} 
        		child { node (a) [person, pathA] {Người A} }
        		child { node (s1) [person] {Em của A} }
        	}
        	child { node (b) [person, pathB] {Người B} };
        	
        	\draw[line, red, ultra thick] (a) -- (c1);
        	\draw[line, red, ultra thick] (c1) -- (root);
        	\draw[line, green!60!black, ultra thick] (b) -- (root);
        	
        \end{tikzpicture}
    \end{minipage}      
    \caption{Mã giả và hình vẽ minh hoạ cho pha 1 của thuật toán.}
    \label{fig:LCAPhase1}
\end{figure}

\vspace*{-\baselineskip}
\textbf{\textit{Pha thứ hai: Xác định Tổ tiên chung gần nhất (LCA).}}
\begin{itemize}
    \item Mục tiêu của pha: Tìm ra nút chung đầu tiên mà cả hai đối tượng $A$ và $B$ cùng thuộc về nhánh hậu duệ của nút đó.
    \item Phương pháp thực hiện: Ta thực hiện phép so sánh hai tập $P_A$ và $P_B$ đã có được từ pha thứ nhất, cụ thể:
    \begin{itemize}
        \item So sánh hai danh sách $P_A$ và $P_B$ từ vị trí nút gốc (đầu danh sách).
        \item Nút cuối cùng trùng khớp trong cả hai danh sách trước khi đường đi rẽ nhánh chính là LCA cần tìm.
    \end{itemize}
    \item Về độ phức tạp: Độ phức tạp thời gian là $O(K)$, với $K = \min(h_A, h_B)$, phép so sánh được thực hiện trực tiếp trên địa chỉ bộ nhớ thay vì so sánh chuỗi.
\end{itemize}

\begin{figure}[H]
	\vspace*{-\baselineskip}
    \centering
    \begin{minipage}{0.50\textwidth}
        \begin{algorithm}[H]
        	\footnotesize
            \caption{DetermineRelationship (A, B)} \label{alg:DetermineRelationship}
            \begin{algorithmic}
                \State Đặt PathA = GetPathToRoot(A)
                \State Đặt PathB = GetPathToRoot(B)
                \For{$i$ from $0$ to $\min(h_A,h_B)$} 
                    \If {$PathA_i$ trùng $PathB_i$}
                        \State Đặt $\text{LCAIndex}=i$
                    \Else \ \textbf{break}
                    \EndIf
                \EndFor
                \If{$\text{LCAIndex}=-1$} \State \Return Không cùng huyết thống
                    \Else \ \Return LCAIndex
                \EndIf
            \end{algorithmic}
        \end{algorithm}
    \end{minipage}%
    \hfill
    \begin{minipage}{0.50\textwidth}
        \centering
        \begin{tikzpicture}[level distance=1.5cm, sibling distance=2.5cm]
        	% Chỉ tô đậm LCA
        	\node (root) [person, LCA] {Cụ Tổ (LCA)}
        	child { node (c1) [person] {Cha A} 
        		child { node (a) [person, pathA] {Người A} }
        		child { node (s1) [person] {Em của A} }
        	}
        	child { node (b) [person, pathB] {Người B} };
        	
        	\draw[line] (a) -- (c1);
        	\draw[line] (c1) -- (root);
        	\draw[line] (b) -- (root);
        \end{tikzpicture}
    \end{minipage}
    \caption{Mã giả và hình vẽ minh hoạ cho pha 2 của thuật toán}
    \label{fig:LCAPhase2}
    \vspace*{-\baselineskip}
\end{figure}


\textbf{\textit{Pha thứ ba: Tính toán véc-tơ khoảng cách.}}
\begin{itemize}
	\item Mục tiêu của pha: Định lượng khoảng cách thế hệ giữa các đối tượng và tổ tiên chung để làm cơ sở cho việc phân giải danh xưng.
	\item Phương pháp thực hiện: Gọi $d_A$ là khoảng cách thế hệ từ $LCA$ đến $A$. Gọi $d_B$ là khoảng cách thế hệ từ $LCA$ đến $B$. Lúc này, ta có một cặp tọa độ $(d_A, d_B)$ phản ánh cấp bậc của hai người, làm tiền đề cho pha cuối cùng.
	\item Về độ phức tạp: Vì mảng đường đi chứa tuần tự các đỉnh từ Gốc đến nút đích, khoảng cách $d$ được chứng minh bằng công thức toán học tuyệt đối: $d = (M - 1) - T$ với $M$ là chiều dài mảng và $T$ là vị trí của LCA. Vậy độ phức tạp thời gian và không gian cho pha này chỉ là $O(1)$.
\end{itemize}

\textbf{\textit{Pha cuối cùng: Phân giải danh xưng}}
\begin{itemize}
	\item Mục tiêu của pha: Chuyển đổi các giá trị toán học $(d_A, d_B)$ thành tên gọi quan hệ họ hàng theo quy chuẩn văn hóa.
	\item Phương pháp thực hiện: Sử dụng cấu trúc điều kiện để ánh xạ cặp $(d_A, d_B)$ vào ma trận danh xưng, cụ thể:
	
	\begin{center}
		\small
		\setlength{\tabcolsep}{4pt}
		\renewcommand{\arraystretch}{1.2}
		\begin{longtable}{
				|>{\centering\arraybackslash}m{2.2cm}|
				>{\centering\arraybackslash}m{3.2cm}|
				>{\centering\arraybackslash}m{3.2cm}|
				>{\arraybackslash}m{6.2cm}|
			}
			\hline
			\textbf{Tọa độ} $(dA,dB)$ & 
			\textbf{Quan hệ của B đối với A} & 
			\textbf{Quan hệ của A đối với B} & 
			\textbf{Giải thích Logic} \\ 
			\hline
			\endfirsthead
			
			\hline
			\textbf{Tọa độ} $(dA,dB)$ & 
			\textbf{Quan hệ của B đối với A} & 
			\textbf{Quan hệ của A đối với B} & 
			\textbf{Giải thích Logic} \\ 
			\hline
			\endhead
			
			\hline
			\multicolumn{4}{r}{\textit{Còn tiếp}} \\
			\endfoot
			
			\hline
			\endlastfoot
			
			$(k,0)$      & Tổ tiên đời thứ $k$                & Hậu duệ đời thứ $k$                & $B \equiv L$, $B$ là gốc của nhánh chứa $A$. \\ \hline
			$(0,k)$      & Hậu duệ đời thứ $k$                & Tổ tiên đời thứ $k$                & $A \equiv L$, $A$ là gốc của nhánh chứa $B$. \\ \hline
			$(k,k)$      & Anh/Chị/Em họ đời $k-1$           & Anh/Chị/Em họ đời $k-1$           & Chung tổ tiên cách đây $k$ thế hệ. \\ \hline
			$(1,2)$      & Bác/Cô/Chú/Dì                      & Cháu (gọi bằng Bác/Chú...)         & $B$ là con của anh/chị/em với $A$. \\ \hline
			$(2,1)$      & Cháu (gọi bằng Bác/Chú...)         & Bác/Cô/Chú/Dì                      & $A$ là con của anh/chị/em với $B$. \\ \hline
			$(1,k)$ với $k>2$ & Cụ/Kỵ (nhánh bàng hệ)          & Chắt/Chút (nhánh bàng hệ)          & $B$ là hậu duệ xa của anh/chị/em với $A$. \\ \hline
			$(dA,dB)$    & Họ hàng tổng quát                  & Họ hàng tổng quát                  & Cách nhau tổng $dA+dB$ bậc thế hệ. \\ \hline
			
		\end{longtable}
	\end{center}
\end{itemize}

\vspace*{-2\baselineskip}

\subsubsection{Thuật toán truy vấn theo danh xưng}

Một bài toán khác được đặt ra ngược với bài toán được xác định ở Mục \ref{sssec:lca1}, đó là tìm người theo quan hệ, nghĩa là: cho một người xác định và quan hệ (bố mẹ/ông bà/cô chú/...), ta cần phải tìm được một (hoặc nhiều) người có quan hệ đó với người đã xác định. Hệ thống tận dụng tính đối ngẫu của thuật toán LCA trước đó. Thay vì tìm tọa độ từ hai nút, ta sẽ cố định tọa độ $(d_A, d_B)$ dựa trên danh xưng tiếng Việt và thực hiện truy vết ngược từ con trỏ \texttt{pParent} kết hợp duyệt nhánh ngang \texttt{pNextSibling} để trả về kết quả cần thiết. Cụ thể, thuật toán này trải qua ba pha như sau:

\textbf{\textit{Pha thứ nhất: Xác định tổ tiên.}}
\begin{itemize}
	\item Mục tiêu của pha: Tìm ra nút tổ tiên chung ($L$).
	\item Phương pháp thực hiện: Từ nút $A$, di chuyển qua con trỏ pParent đúng $d_A$ lần. Ví dụ: Tìm ``Chú'' thì $d_A = 2$. Từ $A$, ta cần leo lên ``Bố'' ($d=1$), rồi sau đó lên ``Ông nội'' ($d=2$). Nút Ông nội chính là $L$.
	\item Về độ phức tạp: Độ phức tạp thời gian là $O(d_A)$, vì chỉ đi theo một đường thẳng nên rất nhanh.
\end{itemize}

\begin{figure}[htbp]
	\vspace*{-\baselineskip}
	\begin{algorithm}[H]
		\footnotesize
		\caption{FindAncestor(startP, dA)}
		\begin{algorithmic}[1]
			\State temp $\gets$ startP
			\For{$i=1$ \to $dA$}
			\If{temp.parent $\neq$ NULL}
			\State temp $\gets$ temp.parent
			\Else \ \Return NULL \Comment{Không đủ độ sâu thế hệ}
			\EndIf
			\EndFor
			\State\ \Return temp
		\end{algorithmic}
	\end{algorithm}
	\vspace*{-\baselineskip}
	\caption{Mã giả mô tả pha thứ nhất của thuật toán truy vấn theo danh xưng}
	\vspace*{-\baselineskip}
\end{figure}


\textbf{\textit{Pha thứ hai: Duyệt và lọc đối tượng.}}
\begin{itemize}
	\item Mục tiêu của pha: Thu thập được những người phù hợp.
	\item Phương pháp thực hiện: Từ tổ tiên $L$, ta cần tìm các hậu duệ ở mức thế hệ $d_B$ thỏa mãn các điều kiện về giới tính và thứ tự. Ta sẽ sử dụng phương pháp duyệt theo chiều rộng (BFS) để tìm tất cả các nút cách $L$ đúng $d_B$ bậc.
	\item Về độ phức tạp: Độ phức tạp thời gian là $O(S)$, với $S$ là số lượng anh/em/con/cháu trong nhánh đang xét đó.
\end{itemize}
%Mã giả mô tả thuật toán của pha ở hình \ref{fig:Rel1}.
\begin{figure}[htbp]
	\vspace*{-\baselineskip}
	\centering
	\begin{algorithm}[H]
		\footnotesize
		\caption{Thuật toán thu thập ứng viên cùng thế hệ}
		\begin{algorithmic}[1]
			% Hàm thu thập tất cả các node cách gốc một khoảng dB
			\Function{CollectCandidates}{root, dB}
			\State Cands $\gets \emptyset$ \Comment{Khởi tạo tập ứng viên rỗng}
			
			\If{$dB = 0$}
			\State \Return \{root\} \Comment{Nếu khoảng cách = 0, trả về chính gốc}
			\EndIf
			
			% Thủ tục đệ quy để duyệt cây
			\Procedure{DuyetTang}{curr, depth}
			\If{curr = NULL} 
			\State \Return \Comment{Trường hợp cơ sở: không có node để xử lý}
			\EndIf
			
			\If{depth = dB}
			\State Cands $\gets$ Cands $\cup$ \{curr\} \Comment{Thêm node hiện tại vào tập ứng viên}
			\State \Return \Comment{Dừng duyệt sâu hơn}
			\EndIf
			
			\State \Call{DuyetTang}{curr.firstChild, depth + 1} \Comment{Duyệt node con đầu tiên (xuống sâu hơn)}
			
			\State t $\gets$ curr.firstChild \Comment{Bắt đầu duyệt các node anh em}
			\While{t $\neq$ NULL \textbf{and} t.nextSibling $\neq$ NULL}
			\State t $\gets$ t.nextSibling \Comment{Chuyển sang node anh em kế tiếp}
			\State \Call{DuyetTang}{$t$, $depth + 1$} \Comment{Duyệt cây con của node anh em này}
			\EndWhile
			\EndProcedure
			
			\State \Call{DuyetTang}{root, 0} \Comment{Bắt đầu duyệt từ gốc ở độ sâu 0}
			\State \Return $Cands$ \Comment{Trả về tất cả các ứng viên đã thu thập được}
			\EndFunction
		\end{algorithmic}
	\end{algorithm}
	\vspace*{-\baselineskip}
	\caption{Mã giả mô tả pha thứ hai của thuật toán truy vấn theo danh xưng}
	\label{fig:Rel1}
\end{figure}
\textbf{\textit{Pha thứ ba: Phân giải danh xưng cụ thể.}}
\begin{itemize}
	\item Mục tiêu của pha: Áp dụng những ràng buộc đã có để đưa ra người cần tìm kiếm.
	\item Phương pháp thực hiện: Đây là lúc áp dụng các ràng buộc mà nhóm đã thiết kế trong struct \texttt{Person}, cụ thể là phân giải theo ``Giới tính'' và ``Thứ tự''. Cụ thể:
	\begin{itemize}
		\item Theo Giới tính: Nếu yêu cầu là ``Cô'', lọc những người có \texttt{gender} là ``Nữ''.
		\item Theo Thứ tự: Dựa vào con trỏ \texttt{pNextSibling}. Nếu người đó nằm phía trước Cha của A trong danh sách anh em, đó là ``Bác'', ngược lại là ``Chú''.
	\end{itemize}
\end{itemize}

Tổng quát, thuật toán chung của ta vẫn tiệm cận mức $O(H)$, rõ ràng là hiệu quả hơn nhiều so với việc quét toàn bộ người trong gia phả. Ngoài ra, thuật toán này còn có ưu điểm là xử lý được trường hợp đa kết quả (ví dụ: tìm ``Chú'' có thể cho ra nhiều người, việc tìm người cụ thể do người dùng quyết định).

Ngoài ra, vì cấu trúc LCRS truyền thống chỉ lưu trữ các thực thể có quan hệ huyết thống trực tiếp, khiến các thành viên nữ (vợ) trở thành ``điểm mù'' (nghĩa là, không thể thực hiện bất kỳ chức năng nào của hệ thống trên những người này) của giải thuật duyệt cây. Vì vậy, giải pháp cho điều này là sử dụng cơ chế Proxy Node (Nút đại diện). Khi truy vấn ``Mẹ'', hệ thống thực hiện tìm kiếm gián tiếp: \texttt{Mother(A) = Spouse(Parent(A))}, cụ thể:
\begin{itemize}
	\item Hệ thống xác định nút Cha của A (là nút $P$ trong cây).
	\item Sau đó, hệ thống truy xuất vào thuộc tính \texttt{spouseName} (vợ) của nút $P$.
\end{itemize}

Cách làm này giữ cho cây luôn là cây đơn gốc (single-root tree), giúp các thuật toán duyệt đệ quy và LCA luôn đạt độ phức tạp tối ưu nhất.

\subsubsection{Thuật toán lưu trữ và phục hồi}

Trong môi trường quản lý dữ liệu bằng con trỏ động, địa chỉ bộ nhớ của các thực thể (\texttt{Person}) chỉ có giá trị nội hàm trong phiên làm việc hiện tại của RAM. Để đảm bảo tính bền vững của dữ liệu, nhóm áp dụng kỹ thuật Tuần tự hóa nhị phân (Binary Serialization) kết hợp với thuật toán Phẳng hóa đồ thị cây.

\textbf{\textit{Cấu trúc Bản ghi tĩnh trung gian (\texttt{PersonRecord})}}

Để khắc phục nhược điểm của chuỗi động \texttt{std::string}, chương trình định nghĩa một cấu trúc \texttt{PersonRecord} kích thước cố định, sử dụng mảng ký tự tĩnh (\texttt{char[]}). Các liên kết con trỏ LCRS (\texttt{pParent}, \texttt{pFirstChild}, \texttt{pNextSibling}) được thay thế hoàn toàn bằng định danh số nguyên là \texttt{id} (mã cá nhân) và \texttt{parentId} (mã của nút cha).

\textbf{\textit{Giải thuật Ghi tệp (Phẳng hóa cây LCRS)}}
\begin{itemize}
	\item Hệ thống sử dụng thuật toán duyệt cây Tiền thứ tự (Pre-order Traversal) thông qua hàm đệ quy \texttt{FlattenTree}.
	\item Tại mỗi nút, dữ liệu được sao chép an toàn (sử dụng \texttt{strncpy} để chống tràn bộ đệm) sang cấu trúc \texttt{PersonRecord} và đẩy vào một mảng động \texttt{std::vector}.
	\item Thay vì ghi tuần tự từng cá nhân, chương trình áp dụng kỹ thuật Block I/O: tính toán tổng dung lượng của mảng và sử dụng phương thức \texttt{write()} để đổ toàn bộ khối dữ liệu từ RAM xuống tệp \texttt{.dat} trong một chu kỳ ghi duy nhất.
\end{itemize}

\begin{figure}[htbp]
	\vspace*{-\baselineskip}
	\centering
	\begin{algorithm}[H]
		\footnotesize
		\caption{Hàm FlattenTree}
		\begin{algorithmic}[1]
			\Require currentNode, recordArray
			
			\State \Comment{Điều kiện dừng đệ quy}
			\If{currentNode là RỖNG}
			\Return
			\EndIf
			
			\State \Comment{1. Chuyển đổi Node động thành Bản ghi tĩnh}
			\State Tạo NewRecord
			\State NewRecord.id = currentNode.id
			
			\If{currentNode có Cha}
			\State NewRecord.parentId = currentNode.parent.id
			\Else
			\State NewRecord.parentId = 0  \Comment{0 ám chỉ đây là Cụ Tổ (Gốc)}
			\EndIf
			
			\State Chép các thông tin (Tên, Năm sinh...) từ Node sang Bản ghi
			
			\State \Comment{2. Lưu vào mảng}
			\State Thêm NewRecord vào recordArray
			
			\State \Comment{3. Đệ quy đi tiếp xuống nhánh Con và nhánh Em}
			\State FlattenTree(currentNode.firstChild, recordArray)
			\State FlattenTree(currentNode.nextSibling, recordArray)
		\end{algorithmic}
	\end{algorithm}
	\vspace*{-\baselineskip}
	\caption{Mã giả mô tả thuật toán phẳng hoá cây để lưu file}
	\vspace*{-\baselineskip}
\end{figure}


\textbf{\textit{Giải thuật Đọc tệp và Tái cấu trúc}}

Quá trình phục hồi dữ liệu từ tệp nhị phân được tối ưu hóa ở mức cao nhất nhằm tránh việc tìm kiếm lặp lại:
\begin{enumerate}
	\item Chương trình đọc toàn bộ khối dữ liệu từ tệp lên một mảng \texttt{PersonRecord} trên RAM.
	\item Khởi tạo một Bảng băm (Hash Map) sử dụng cấu trúc \texttt{std::unordered\_map<int, Person*>} để thiết lập quan hệ ánh xạ giữa \texttt{id} nguyên thủy và địa chỉ bộ nhớ động vừa được cấp phát.
	\item Khi duyệt qua từng bản ghi, hệ thống tạo mới một đối tượng \texttt{Person} và sử dụng \texttt{parentId} để tra cứu trực tiếp địa chỉ người cha trong Bảng băm. Thao tác tra cứu này đạt độ phức tạp thời gian $O(1)$, cho phép liên kết con trỏ ngay lập tức thông qua hàm \texttt{AddChild} mà không cần duyệt lại cây LCRS.
	\item Hệ thống tự động cập nhật lại biến đếm ID toàn cục (\texttt{global\_id\_counter}) dựa trên mức ID lớn nhất trong tệp để đảm bảo tính duy nhất cho các phiên nhập liệu tiếp theo.
\end{enumerate}

\textbf{\textit{Đánh giá thuật toán}}

Cả hai tiến trình Serialization và Deserialization đều đạt độ phức tạp thời gian $O(N)$ và không gian $O(N)$, với $N$ là tổng số thành viên trong gia phả. Việc kết hợp Block I/O và Hash Map giúp giải thuật xử lý mượt mà dữ liệu quy mô hàng vạn thực thể mà không gây ra tình trạng tràn ngăn xếp (Stack Overflow) hay nghẽn cổ chai truy xuất đĩa.

\subsubsection{Thuật toán thống kê dòng họ}

Thuật toán thống kê cung cấp cái nhìn tổng quan về quy mô và cấu trúc của dòng họ thông qua các chỉ số định lượng, nhằm giúp người quản trị nắm bắt được tổng số thành viên (bao gồm cả huyệt thống và phu nhân), sự phân bổ giới tính, độ dài thế hệ (số đời), và tình trạng sinh tồn của các cá nhân trong gia phả.

Do gia phả được lưu trữ dưới dạng cây LCRS, việc thống kê được thực hiện tối ưu nhất bằng phương pháp Duyệt cây theo chiều sâu (DFS) kết hợp với kỹ thuật đệ quy. Mỗi khi đi qua một nút ($Node$), hệ thống sẽ thực hiện kiểm tra các thuộc tính của thực thể Person và cập nhật vào cấu trúc dữ liệu \texttt{FamilyStats}.

Các bước cụ thể như sau:
\begin{itemize}
	\item Bước 1: Tạo một cấu trúc lưu trữ \texttt{FamilyStats} với các biến đếm (tổng số, nam, nữ, phu nhân, tử vong) và một mảng \texttt{genDistribution} để lưu số lượng thành viên của từng đời.
	\item Bước 2: Đệ quy xử lý nút hiện tại. Tăng biến \texttt{totalMembers}; Kiểm tra gender: Nếu là ``Nam'' tăng \texttt{maleCount}, nếu ``Nu'' tăng \texttt{femaleCount}; kiểm tra \texttt{spouseName}: Nếu khác ``None'' hoặc ``N/A'', tăng \texttt{totalSpouses} (thống kê dâu/rể); Kiểm tra \texttt{deathDay}: Nếu khác ``N/A'', tăng \texttt{deceasedCount}; Cập nhật thế hệ: Sử dụng tham số \texttt{level} để tăng giá trị tại \texttt{genDistribution[level]} và cập nhật \texttt{maxGen = max(maxGen, level)}.
	\item Bước 3: Lan toả đệ quy. Gọi đệ quy xuống con đầu tiên và gọi đệ quy sang anh em kế cận.
\end{itemize}

\textbf{\textit{Phân tích độ phức tạp}}

\begin{itemize}
	\item Độ phức tạp thời gian: $O(N)$, với $N$ là tổng số nút trong cây gia phả. Thuật toán chỉ đi qua mỗi nút đúng một lần duy nhất.
	\item Độ phức tạp không gian: $O(H)$, với $H$ là chiều cao của cây (tương ứng với số đời), do sử dụng ngăn xếp đệ quy trong quá trình duyệt.
\end{itemize}